\documentclass[../../../include/open-logic-section]{subfiles}

\begin{document}

\olfileid{sth}{ordinals}{vn}
\olsection[फोन न्यूमनची रचना]{क्रमसूचक संख्यांची फोन न्यूमनची रचना}

\olref[sth][ordinals][iso]{thm:woalwayscomparable} वरून एक कल्पना सुचते. सुक्रमणांचे प्रतिनिधी म्हणून \emph{क्रमप्रकार} नावाच्या काही वस्तू आपण मांडू शकतो. $\tuple{A, <}$ या सुक्रमणाचा क्रमप्रकार $\ordtype{A, <}$ असा लिहू; मग पुढील दोन तत्त्वे मिळतील अशी अपेक्षा असेल:
\begin{align*}
	\ordtype{A, <} = \ordtype{B, \lessdot} &
	\text{ तेव्हा आणि केवळ तेव्हाच } \ordeq{\tuple{A, <}}{\tuple{B, \lessdot}}\\
	\ordtype{A, <} < \ordtype{B, \lessdot}&
	\text{ तेव्हा आणि केवळ तेव्हाच }\ordeq{\tuple{A, <}}{\tuple{B_b, \lessdot_b}}\text{ असे काही }b \in B
\end{align*}
शिवाय, नैसर्गिक संख्या जशा विशिष्ट संचांच्या रूपात मांडता येतात, तसे क्रमप्रकारही \emph{विशिष्ट संचांच्या रूपात} मांडता येतील अशी अपेक्षा असू शकते.

ते करण्याची सर्वाधिक प्रचलित पद्धत---आणि आपण अनुसरणार असलेली पद्धत---म्हणजे विशिष्ट \emph{प्रमाणभूत} सुक्रमित संचांच्या आधारे हे क्रमप्रकार परिभाषित करणे. असे प्रमाणभूत संच प्रथम फोन न्यूमन यांनी मांडले:

\begin{defn}
$(\forall x \in A)x \subseteq A$ असेल तेव्हा आणि केवळ तेव्हाच $A$ हा संच \emph{संक्रमक} असतो. $A$ हा संक्रमक असून त्याच्या सदस्यांपुरत्या $\in$ च्या बंधनाने सुक्रमित असेल तेव्हा आणि केवळ तेव्हाच $A$ ही \emph{क्रमसूचक संख्या} असते.
\end{defn}
\noindent
यापुढे क्रमसूचक संख्यांसाठी ग्रीक अक्षरे वापरू. व्याख्येवरून लगेच दिसते की $\alpha$ ही क्रमसूचक संख्या असल्यास $\tuple{\alpha, \in_\alpha}$ हे सुक्रमण आहे; येथे $\in_\alpha =
\Setabs{\tuple{x, y} \in \alpha^2}{x \in y}$. म्हणून संकेतलेखनात किंचित मोकळीक घेऊन $\alpha$ \emph{स्वतःच} एक सुक्रमण आहे असेही म्हणू शकतो.

पहिल्या काही क्रमसूचक संख्या पुढीलप्रमाणे आहेत:
\[
	\emptyset, \{\emptyset\},
	\{\emptyset, \{\emptyset\}\},
	\{\emptyset, \{\emptyset\}, \{\emptyset, \{\emptyset\}\}\}, \ldots
\]
अनंततेच्या स्वयंसिद्धकात, म्हणजे फोन न्यूमन यांच्या $\omega$ च्या व्याख्येत, आपल्याला भेटलेल्या पहिल्या काही क्रमसूचक संख्या याच आहेत (\olref[sth][z][infinity-again]{sec} पाहा). हा योगायोग नाही. फोन न्यूमन यांच्या व्याख्येत नैसर्गिक संख्या क्रमसूचक संख्या असतात; पण अनंतपार क्रमसूचक संख्यांनाही जागा असते.

नेहमीप्रमाणे आता विचारता येईल: \emph{याच} क्रमसूचक संख्या आहेत का? की फोन न्यूमन यांनी आपल्याला असे काही संच दिले आहेत ज्यांना आपण क्रमसूचक संख्या \emph{मानू शकतो}? या प्रश्नावरील चर्चा \olref[sfr][rel][ref]{sec}, \olref[sfr][arith][ref]{sec}, \olref[sfr][infinite][dedekindsproof]{sec} आणि \olref[sth][z][nat]{sec} मधील चर्चांप्रमाणेच आहे; म्हणून येथे ती पुन्हा विस्ताराने करणार नाही. पुढे ``क्रमसूचक संख्या'' म्हटल्यावर आपण ``फोन न्यूमन यांच्या क्रमसूचक संख्या'' अभिप्रेत धरू.

\end{document}
